\chapter{Introduction}
\label{chap:introduction}

\researchstatus{Early authorial draft}{This chapter preserves the moral
purpose and opening intuitions of the project. Its claims are invitations to
research, not universal cryptographic guarantees.}

\section{Author's Intuition}

Assume that an adversary may be smarter than us. The original intuition hoped
that greater understanding might also yield kindness. That hope motivates the
book, but no protocol may assume that intelligence implies benevolence.

Here an \emph{adversary} is a party willing to pursue a shared purpose while
protecting its own interests. An \emph{enemy} is provisionally distinguished
as a party that acts against the shared purpose itself. Later chapters must
replace this moral vocabulary with explicit roles, actions, and adversarial
capabilities whenever a technical claim depends on it.

\section{The Project's Promise}

One generation has often been told that war will last forever. Differences and
resistance can nevertheless preserve hope under oppression. This project asks
whether some conflicts contain procedural weaknesses: places where parties
could commit, prove, compute, or coordinate without surrendering every secret
or trusting a single intermediary.

The book does not claim that every unnecessary war has already been solved.
It studies smaller questions that can be falsified. For each proposed protocol
we ask what parties know, what they keep private, what they can verify, how
they may deviate, and how the protocol ends.

\section{Humility Before Security Claims}

The early draft imagined formulations that no quantum computer could break and
described them broadly as information-theoretically secure. That statement is
not a property of the book as a whole. Some constructions may seek
information-theoretic guarantees; others rely on classical computational
assumptions; many current chapters are only educational demonstrations or open
problems. Each chapter must state its own model and evidence.

An adversary may be more capable than the designer. Humility therefore comes
before every construction. If a reader sees a missing assumption, a stronger
attack, or a better path, the project asks to be corrected.

\section{The Obvious Lying Adversary}

Zero-knowledge proofs, commitments, authentication, and transcript mechanisms
can constrain particular false statements or inconsistent actions under named
assumptions. They do not prevent people from lying in general, establish just
objectives, or make an institution trustworthy by themselves.

\section{The Book Ahead}

The foundations distinguish randomness, identity, commitment, and proof. The
finite-game chapters then place small protocols under a bright light. Later
research notes ask about private joint action, institutions, failed metrics,
aftermath, and repair. The appendices collect the primitive background needed
to assess those ideas.

Life can become more beautiful. For now, let us investigate how to make war
optional. This is the ethical aspiration of the project, not a theorem proved
by the pages that follow.
